\subsection{信息量、条件熵与相关性}

一个事件越罕见，观察到它所获得的信息越多。要求独立事件的信息可加、确定事件的信息为零，
即 $I(pq)=I(p)+I(q)$、$I(1)=0$；再要求连续及随概率递减，得到 $I(p)=-c\ln p$，其中 $c>0$。
例如令 $p=e^{-x}$，则 $J(x)=I(e^{-x})$ 满足连续加法方程 $J(x+y)=J(x)+J(y)$，故 $J(x)=cx$。
取 $c=1$，平均信息量是
\[H(X)=\sum_i p_i I(p_i)=-\sum_i p_i\ln p_i.\]
自然对数给出的信息单位为奈特；若采用以2为底的对数，单位为比特。
公平骰子的熵是 $\ln6$；已知为偶数后的熵是 $\ln3$，所获得的信息是 $\ln2$。

把一个结果进一步细分为 $(i,j)$，写 $p_{ij}=p_i p(j\mid i)$，代入定义：
\begin{align*}
 H(X,Y)&=-\sum_{ij}p_{ij}\ln[p_i p(j\mid i)]\\
 &= -\sum_i p_i\ln p_i
    -\sum_i p_i\sum_jp(j\mid i)\ln p(j\mid i)\\
 &=H(X)+H(Y\mid X).
\end{align*}
这里 $H(Y\mid X)$ 是先知道 $X$ 后剩余不确定性的加权平均。独立时 $p(j\mid i)=q_j$，
因而 $H(X,Y)=H(X)+H(Y)$。

一般情况下不能直接相加。令 $q_{ij}=p_iq_j$，利用 $-\ln x\ge1-x$，有
\[\sum_{p_{ij}>0}p_{ij}\ln\frac{p_{ij}}{q_{ij}}
\ge\sum_{p_{ij}>0}(p_{ij}-q_{ij})
=1-\sum_{p_{ij}>0}q_{ij}\ge0.\]
零概率项以连续极限解释；若 $p_{ij}=0$ 的位置上 $q_{ij}>0$，上述不等式仍成立。
左边等于 $H(X)+H(Y)-H(X,Y)$，即互信息。所以联合熵不大于两边之和，等号当且仅当独立。
相关性会减少额外信息，但协方差为零并不足以证明独立；不能把两者混淆。
例如用独立、等方差的零均值高斯变量 $X,Y$ 模拟随机光场的两个正交分振幅，
与 $x$ 轴夹角 $\theta$ 方向的分振幅为 $Z=X\cos\theta+Y\sin\theta$。
它与 $X$ 的相关系数为 $\cos\theta$：$\theta=\pi/2$ 时独立；锐角时有关联。
对于这个高斯模型，$I(X{:}Z)=-\tfrac12\ln(1-\cos^2\theta)$，可由二维高斯熵与两个边缘熵相减得到。
同方向的连续、无测量噪声变量完全相同，互信息发散；若按有限分辨率离散读数，同方向的互信息才等于该离散变量的熵。
这说明讨论光场相关性时应交代分布模型和测量分辨率，不能把连续与离散结论直接混用。

\subsection{连续分布与最大熵的条件}

连续变量可形式地写微分熵 $h=-\int f\ln f\,\dd x$，但量纲与坐标选择必须明确。
可以先选定参考尺度把变量无量纲化，或用参考测度写相对熵；不能对带单位的密度直接取对数后比较不同单位下的数值。
对固定的无量纲坐标，给定均值 $\mu$ 和方差 $\sigma^2>0$，令
\[g(x)=\frac1{\sqrt{2\pi}\sigma}\exp\!\left[-\frac{(x-\mu)^2}{2\sigma^2}\right].\]
任意具有相同均值、方差的密度 $f$ 满足 $D(f\Vert g)=\int f\ln(f/g)\dd x\ge0$。
因 $\ln g$ 只含常数和 $(x-\mu)^2$，由约束可算出
\[D(f\Vert g)=-h(f)+\frac12\ln(2\pi\sigma^2)+\frac12\ge0.\]
所以 $h(f)\le\tfrac12\ln(2\pi e\sigma^2)$，且仅在 $f=g$ 时取等号。
高斯分布是这些约束下的最大熵分布，不是任何约束下都必然出现的分布。

均匀角分布也须考虑测度：球面上均匀取方向时，面积元为
$\dd A=r^2\sin\theta\,\dd\theta\dd\phi$，故极角密度为 $f_\theta(\theta)=\tfrac12\sin\theta$，
不是在 $0\le\theta\le\pi$ 上均匀。常数由 $\int_0^\pi\sin\theta\dd\theta=2$ 确定。
这是几何概率、变量变换和统计物理中速度空间球壳因子的共同思想。
